%%==========================================================================%%
%% usagenotes.tex -- "Usage Notes". Access/registration pattern modeled on    %%
%% the reference paper, genericized to a credentialed repository (the ref uses %%
%% PublicnEUro/GDPR; we will likely use PhysioNet credentialed access + DUA).  %%
%% REWRITE before submission. [TBD: ...] = fill in.                            %%
%%==========================================================================%%

\section*{Usage Notes}\label{sec:usage}

\subsection*{Access conditions}\label{subsec:license}
% Reusable pattern: state de-identification status -> resulting access model
% (open vs credentialed + DUA), registration steps, expected timelines, license,
% commercial-use terms. Adapt to the target repository's process.
[TBD: state the de-identification status of the release and the resulting access
model.] Access is provided via [TBD: repository, e.g., PhysioNet], where users
[TBD: register / become credentialed and sign a Data Use Agreement] before
download. [TBD: expected timelines, the data license, and any restrictions on
commercial use.]

\subsection*{Accessing the imaging and metadata}\label{subsec:access}
The imaging is distributed as NIfTI (\texttt{.nii.gz}) and can be read with
standard tools [TBD: e.g., \texttt{nibabel}, SimpleITK]. The accompanying
plain-text metadata tables can be loaded with any standard data-analysis tooling.

\subsection*{Querying the index}\label{subsec:query}
% Show how the SQLite index enables cohort selection.
The provided relational index supports cohort selection --- for example,
selecting subjects by age band, sequence/contrast, pathology category, label
availability, or annotating reader. [TBD: include example queries.]

\subsection*{Query/preview client}\label{subsec:client}
% OPEN DECISION: thin-client scope (recommended: optional lightweight client;
% see CLAUDE.md). Dataset must remain fully usable without it.
% TODO (CLAUDE.md): check whether 3D Slicer can filter by sequence (e.g., DWI);
% if so, it may serve as (or complement) the preview client.
A lightweight client is provided for browsing and querying the dataset over the
relational index; the dataset is fully usable without it. [TBD: confirm whether
3D Slicer can filter by sequence (e.g., DWI) and whether it can serve as a
preview path.]

\subsection*{Multimodal use of radiology reports}\label{subsec:multimodal}
Where paired scrubbed radiology reports are included, they can be combined with
the imaging for multimodal modelling --- e.g., image--text pretraining, report
generation, or weak supervision for pathology labels.

\subsection*{Suggested splits, intended uses, and limitations}\label{subsec:limits}
% Recommended split handling (avoid subject leakage), intended tasks, known
% limitations (cohort/site bias, sequence heterogeneity, label scope).
We recommend subject-level data splits to avoid leakage across train/validation/
test partitions. [TBD: state intended tasks and known limitations --- e.g.,
single-site/cohort bias, heterogeneous protocols, and the scope of available
labels.]
